\documentclass[12pt,a4paper]{article}
\usepackage[utf8]{inputenc}
\usepackage[T1]{fontenc}
\usepackage[brazilian]{babel}
\usepackage{amsmath,amssymb}
\usepackage{booktabs,array,longtable}
\usepackage{tikz}
\usepackage[margin=2.5cm]{geometry}
\usepackage{setspace}
\usepackage{enumitem}
\onehalfspacing
\setlength{\parindent}{0pt}
\setlength{\parskip}{0.6em}
\newcommand{\xb}{\bar{x}}

\begin{document}

\begin{center}
{\large Fatec Baixada Santista Rubens Lara}\\
Centro Paula Souza - Governo do Estado de São Paulo\\[1.2em]
{\Large\textbf{1ª Atividade Prática de Estatística Indutiva}}\\[0.6em]
Curso de Ciência de Dados - 3º ciclo\\
Profª. Márcia Roberta S. Pires Silva\\[1.2em]
\textbf{Integrantes}\\
Diego Pacheco Hissnauer\\
Kauan Nunes Vieira Santos\\
Rafael Manfe Bisaio\\
Victor Gabriel Leme da Silva
\end{center}

\vspace{1em}

\section*{Apresentação/Contexto}

O arquivo \texttt{clientes\_plataforma\_Streaming.csv} reúne informações de $N = 100\,000$ clientes de uma plataforma de streaming. A variável estudada é \texttt{Gasto\_reais}, o gasto de cada cliente, em reais.

Neste relatório, os cálculos seguem o método visto em aula: valores críticos lidos nas tabelas impressas (Tabela III, normal padrão; Tabela V, t de Student), variância amostral com divisor $n-1$, graus de liberdade $\nu = n-1$ e conclusões escritas na linguagem do problema. O software R foi usado para sortear as amostras e para conferir as contas; o código está no repositório do trabalho (\texttt{streaming-spend-inference}, no GitHub) e todos os sorteios usam a semente 42, o que permite repetir exatamente os mesmos resultados.

\section{Parte 1: Amostra-piloto}

\subsection*{Como a amostra foi sorteada}

A população é formada pelos $N = 100\,000$ clientes, que já vêm numerados de 1 a $N$ pela coluna \texttt{Cliente\_ID}. Isso nos dá o quadro amostral, a lista de todos os elementos da população, necessária para a amostragem aleatória simples (AAS).

Foram sorteados $n = 25$ números distintos entre 1 e 100\,000, ou seja, uma AAS \textbf{sem reposição}, porque cada cliente deve aparecer no máximo uma vez. Cada cliente tem a mesma probabilidade de entrar na amostra:
$$P = \frac{n}{N} = \frac{25}{100\,000} = 0{,}00025 = 0{,}025\%.$$

No R, o sorteio foi feito com \texttt{set.seed(42)} seguido de \texttt{sample(1:100000, 25)}.

\subsection*{(a) Dados da amostra-piloto}

Seja $X$ o gasto, em reais, de um cliente da plataforma, e $X_1, \dots, X_{25}$ os gastos dos 25 clientes sorteados.

\begin{center}
\begin{tabular}{rrr}
\toprule
$i$ & \texttt{Cliente\_ID} & Gasto $x_i$ (R\$) \\
\midrule
1 & 61413 & 307,21 \\
2 & 54425 & 375,28 \\
3 & 99556 & 1\,405,57 \\
4 & 74362 & 376,63 \\
5 & 46208 & 1\,467,55 \\
6 & 47128 & 1\,262,81 \\
7 & 16740 & 707,08 \\
8 & 61605 & 329,86 \\
9 & 73236 & 905,18 \\
10 & 9091 & 589,42 \\
11 & 62041 & 1\,498,96 \\
12 & 33700 & 521,75 \\
13 & 59359 & 814,28 \\
14 & 54789 & 442,49 \\
15 & 54586 & 689,20 \\
16 & 13610 & 455,13 \\
17 & 28559 & 1\,187,84 \\
18 & 88261 & 1\,185,36 \\
19 & 11224 & 995,16 \\
20 & 33713 & 709,38 \\
21 & 64410 & 175,72 \\
22 & 77954 & 207,86 \\
23 & 81301 & 583,06 \\
24 & 74743 & 568,56 \\
25 & 96954 & 1\,387,04 \\
\bottomrule
\end{tabular}
\end{center}

\subsection*{(b) Média e variância}

A média amostral é
$$\xb = \frac{1}{n}\sum_{i=1}^{n} x_i = \frac{19\,148{,}38}{25} = 765{,}9352.$$

A variância amostral usa o divisor $n-1$ (aula 03), porque $\xb$ já foi calculada com os mesmos dados e, por isso, apenas $n-1$ desvios podem variar livremente:
$$S^2 = \frac{1}{n-1}\sum_{i=1}^{n}(x_i - \xb)^2.$$

Para fazer a soma, calculamos primeiro o desvio de cada valor em relação à média e o quadrado desse desvio:

{\footnotesize\renewcommand{\arraystretch}{0.92}
\begin{longtable}{rrrr}
\toprule
$i$ & $x_i$ & $x_i - \xb$ & $(x_i - \xb)^2$ \\
\midrule
\endhead
1 & 307,21 & $-$458,7252 & 210\,428,8091 \\
2 & 375,28 & $-$390,6552 & 152\,611,4853 \\
3 & 1\,405,57 & 639,6348 & 409\,132,6774 \\
4 & 376,63 & $-$389,3052 & 151\,558,5387 \\
5 & 1\,467,55 & 701,6148 & 492\,263,3276 \\
6 & 1\,262,81 & 496,8748 & 246\,884,5669 \\
7 & 707,08 & $-$58,8552 & 3\,463,9346 \\
8 & 329,86 & $-$436,0752 & 190\,161,5801 \\
9 & 905,18 & 139,2448 & 19\,389,1143 \\
10 & 589,42 & $-$176,5152 & 31\,157,6158 \\
11 & 1\,498,96 & 733,0248 & 537\,325,3574 \\
12 & 521,75 & $-$244,1852 & 59\,626,4119 \\
13 & 814,28 & 48,3448 & 2\,337,2197 \\
14 & 442,49 & $-$323,4452 & 104\,616,7974 \\
15 & 689,20 & $-$76,7352 & 5\,888,2909 \\
16 & 455,13 & $-$310,8052 & 96\,599,8723 \\
17 & 1\,187,84 & 421,9048 & 178\,003,6603 \\
18 & 1\,185,36 & 419,4248 & 175\,917,1629 \\
19 & 995,16 & 229,2248 & 52\,544,0089 \\
20 & 709,38 & $-$56,5552 & 3\,198,4906 \\
21 & 175,72 & $-$590,2152 & 348\,353,9823 \\
22 & 207,86 & $-$558,0752 & 311\,447,9289 \\
23 & 583,06 & $-$182,8752 & 33\,443,3388 \\
24 & 568,56 & $-$197,3752 & 38\,956,9696 \\
25 & 1\,387,04 & 621,1048 & 385\,771,1726 \\
\midrule
\multicolumn{3}{r}{Soma dos quadrados dos desvios} & 4\,241\,082,31 \\
\bottomrule
\end{longtable}
}

Os desvios somam zero (a menos de arredondamento), como deve ocorrer em torno da média. Substituindo:
$$S^2 = \frac{4\,241\,082{,}31}{25-1} = \frac{4\,241\,082{,}31}{24} = 176\,711{,}76 \ (\text{R\$}^2).$$

O desvio padrão amostral é $S = \sqrt{S^2} = \sqrt{176\,711{,}76} = 420{,}37$ reais.

\textbf{Resultado da Parte 1:} $\xb = \text{R\$ } 765{,}9352$ \ e \ $S^2 = 176\,711{,}76\ \text{R\$}^2$.

\section{Parte 2: Intervalo de confiança}

\subsection*{Escolha da distribuição}

Seja $X$ o gasto de um cliente e $X$ amostral o gasto dos 25 clientes da piloto. O enunciado manda considerar a variância populacional $\sigma^2$ \textbf{desconhecida}. Nesse caso, o desvio padrão da população é substituído por $S$ e a variável padronizada não segue a normal, e sim a distribuição \textbf{t de Student com $\nu = n-1$ graus de liberdade}. Admite-se, como na disciplina, que a média amostral tem comportamento aproximadamente normal (Teorema do Limite Central).

\subsection*{Dados}

$n = 25$, \ $\nu = 24$, \ $\xb = 765{,}9352$, \ $S = 420{,}3710$, \ confiança $1-\alpha = 95\%$, logo $\alpha = 0{,}05$.

\subsection*{Fórmula}

$$IC(\mu;\,95\%) = \left[\ \xb - t_{\alpha/2}\cdot\frac{S}{\sqrt{n}}\ ;\ \xb + t_{\alpha/2}\cdot\frac{S}{\sqrt{n}}\ \right].$$

\subsection*{Cálculos}

\textbf{1. Erro padrão estimado:}
$$\frac{S}{\sqrt{n}} = \frac{420{,}3710}{\sqrt{25}} = \frac{420{,}3710}{5} = 84{,}0742.$$

\textbf{2. Valor crítico.} A Tabela V é bicaudal: a coluna indica a soma das áreas das duas caudas. Para confiança de 95\%, usa-se a coluna $p = \alpha = 5\%$, na linha $\nu = 24$:
$$t_{\alpha/2} = 2{,}064.$$

\textbf{3. Margem de erro:}
$$e = 2{,}064 \cdot 84{,}0742 = 173{,}5291.$$

\textbf{4. Limites do intervalo:}
$$\xb - e = 765{,}9352 - 173{,}5291 = 592{,}4061, \qquad \xb + e = 765{,}9352 + 173{,}5291 = 939{,}4643.$$

$$\boxed{IC(\mu;\,95\%) = [\,592{,}41 \ ;\ 939{,}46\,]}$$

\subsection*{Representação gráfica}

A figura mostra a curva t de Student com 24 graus de liberdade. A área central, de 95\%, fica entre $-2{,}064$ e $+2{,}064$; sobram $2{,}5\%$ em cada cauda, somando $\alpha = 5\%$. O valor $t_{\alpha/2} = 2{,}064$ é o ponto da tabela que separa essas áreas. A margem de erro do intervalo é esse valor multiplicado pelo erro padrão estimado, e foi isso que somamos e subtraímos de $\xb$.

\begin{center}
\begin{tikzpicture}[x=1.3cm,y=8cm]
  \fill[gray!35] (-4,0) -- plot[domain=-4:-2.064,samples=40] (\x,{0.39481*(1+\x*\x/24)^(-12.5)}) -- (-2.064,0) -- cycle;
  \fill[gray!35] (2.064,0) -- plot[domain=2.064:4,samples=40] (\x,{0.39481*(1+\x*\x/24)^(-12.5)}) -- (4,0) -- cycle;
  \fill[blue!10] (-2.064,0) -- plot[domain=-2.064:2.064,samples=60] (\x,{0.39481*(1+\x*\x/24)^(-12.5)}) -- (2.064,0) -- cycle;
  \draw[thick] plot[domain=-4:4,samples=80] (\x,{0.39481*(1+\x*\x/24)^(-12.5)});
  \draw[->] (-4.3,0) -- (4.5,0) node[right] {$t$};
  \draw (-2.064,0) -- (-2.064,-0.012) node[below] {$-2{,}064$};
  \draw (0,0) -- (0,-0.012) node[below] {$0$};
  \draw (2.064,0) -- (2.064,-0.012) node[below] {$+2{,}064$};
  \draw[dashed] (-2.064,0) -- (-2.064,0.0545);
  \draw[dashed] (2.064,0) -- (2.064,0.0545);
  \node at (0,0.17) {\small $1-\alpha = 95\%$};
  \node[font=\footnotesize,align=center] at (-3.25,0.115) {$\alpha/2$\\$=2{,}5\%$};
  \node[font=\footnotesize,align=center] at (3.25,0.115) {$\alpha/2$\\$=2{,}5\%$};
\end{tikzpicture}
\end{center}

\subsection*{Interpretação}

Com 95\% de confiança, o gasto médio dos clientes da plataforma está entre R\$ 592{,}41 e R\$ 939,46. O nível de confiança se refere ao método: se repetíssemos o sorteio da amostra e a construção do intervalo muitas vezes, cerca de 95\% dos intervalos obtidos conteriam a verdadeira média populacional $\mu$.

\section{Parte 3: Determinação do tamanho da amostra}

\subsection*{Ideia}

Na Parte 2, o tamanho da amostra era fixo e calculamos o intervalo. Agora fazemos o caminho inverso: fixamos o \textbf{erro máximo} que aceitamos cometer, $\varepsilon$, e o nível de confiança, e perguntamos qual o menor $n$ que garante isso. Exige-se
$$P\big(|\bar{X} - \mu| \le \varepsilon\big) \ge 1-\alpha.$$

\subsection*{Fórmula}

Como $\bar{X} \sim N(\mu,\ \sigma^2/n)$, padronizamos dividindo pelo erro padrão $\sigma/\sqrt{n}$:
$$P\left(-\frac{\sqrt{n}\,\varepsilon}{\sigma} \le Z \le \frac{\sqrt{n}\,\varepsilon}{\sigma}\right) = 1-\alpha .$$
A área central $1-\alpha$ fica entre $-z$ e $+z$, então o limite superior precisa ser o próprio $z$: $\sqrt{n}\,\varepsilon/\sigma = z$. Elevando ao quadrado e isolando $n$:
$$\boxed{\,n = \frac{z^2\,\sigma^2}{\varepsilon^2}\,}$$

Como $\sigma^2$ é desconhecida, o enunciado manda usar a variância da amostra-piloto no lugar dela: $\sigma^2 \approx S^2$. Usa-se o $z$ da normal, como na aula, porque o $n$ é justamente o que ainda não conhecemos; como o resultado será bem maior que 25, a diferença para a distribuição t é desprezível.

\subsection*{Dados}

$S^2 = 176\,711{,}76$ (Parte 1); \ $\varepsilon = \text{R\$ } 50{,}00$, na mesma unidade da variável (reais); \ confiança $1-\alpha = 95\%$.

\subsection*{Valor crítico}

Na Tabela III, que dá a área entre $0$ e $z$, procuramos a área $(1-\alpha)/2 = 0{,}95/2 = 0{,}47500$. Ela está na linha $1{,}9$, coluna $6$, logo
$$z = 1{,}96, \qquad z^2 = 3{,}8416.$$

\subsection*{Representação gráfica}

A curva é a normal padrão. A área central de 95\% vai de $-1{,}96$ a $+1{,}96$. Na escala da média amostral, esse trecho corresponde a $\bar{X}$ ficar a no máximo $\varepsilon$ de $\mu$. Escolher $n$ é ajustar o erro padrão $\sigma/\sqrt{n}$ para que $1{,}96$ erros padrão caibam exatamente em $\varepsilon = 50$.

\begin{center}
\begin{tikzpicture}[x=1.3cm,y=8cm]
  \fill[blue!10] (-1.96,0) -- plot[domain=-1.96:1.96,samples=60] (\x,{0.39894*exp(-\x*\x/2)}) -- (1.96,0) -- cycle;
  \fill[gray!35] (-4,0) -- plot[domain=-4:-1.96,samples=40] (\x,{0.39894*exp(-\x*\x/2)}) -- (-1.96,0) -- cycle;
  \fill[gray!35] (1.96,0) -- plot[domain=1.96:4,samples=40] (\x,{0.39894*exp(-\x*\x/2)}) -- (4,0) -- cycle;
  \draw[thick] plot[domain=-4:4,samples=80] (\x,{0.39894*exp(-\x*\x/2)});
  \draw[->] (-4.3,0) -- (4.5,0) node[right] {$z$};
  \draw (-1.96,0) -- (-1.96,-0.012) node[below] {$-1{,}96$};
  \draw (0,0) -- (0,-0.012) node[below] {$0$};
  \draw (1.96,0) -- (1.96,-0.012) node[below] {$+1{,}96$};
  \draw[dashed] (-1.96,0) -- (-1.96,0.0585);
  \draw[dashed] (1.96,0) -- (1.96,0.0585);
  \node[font=\footnotesize] at (-1.96,-0.098) {$\bar{x}-\mu = -\varepsilon$};
  \node[font=\footnotesize] at (0,-0.098) {$\bar{x}-\mu = 0$};
  \node[font=\footnotesize] at (1.96,-0.098) {$\bar{x}-\mu = +\varepsilon$};
  \node at (0,0.17) {\small $1-\alpha = 95\%$};
\end{tikzpicture}
\end{center}

\vspace{1.2em}

\subsection*{Cálculos}

$$n = \frac{z^2\,S^2}{\varepsilon^2} = \frac{3{,}8416 \cdot 176\,711{,}76}{50^2} = \frac{678\,855{,}91}{2\,500} = 271{,}5424.$$

O tamanho da amostra é sempre arredondado \textbf{para cima}: com 271 clientes o erro ficaria um pouco acima de 50, e arredondar para baixo não cumpriria a exigência. Logo,
$$\boxed{n_{\text{mín}} = 272 \text{ clientes}}$$

\subsection*{Verificação}

A margem de erro obtida com cada tamanho é $z\cdot S/\sqrt{n}$, com $z\cdot S = 1{,}96 \cdot 420{,}3710 = 823{,}9271$:
\begin{align*}
n = 271:&\quad \frac{823{,}9271}{\sqrt{271}} = \frac{823{,}9271}{16{,}4621} = 50{,}05 \quad (\text{maior que } 50),\\[0.4em]
n = 272:&\quad \frac{823{,}9271}{\sqrt{272}} = \frac{823{,}9271}{16{,}4924} = 49{,}96 \quad (\text{menor que } 50).
\end{align*}

\subsection*{Conclusão}

Para estimar o gasto médio dos clientes com 95\% de confiança e erro máximo de R\$ 50,00, são necessários \textbf{no mínimo 272 clientes} na amostra. Esse valor depende da variância estimada na amostra-piloto: outra piloto daria outra estimativa de $\sigma^2$ e, portanto, outro $n$.

\section{Parte 4: Distribuição da média amostral}

\textit{Em elaboração.}

\section{Parte 5: Teste de hipótese para a média}

\textit{Em elaboração.}

\end{document}
